\section{Test Case 2: 1-D case with bead insertion at the nozzle activated}

This input file (located in the \textit{input-2} sub-directory) keeps the
material parameters, the jet radius at the nozzle, the collector distance and
the external potential of Test Case 1, but starts from an initial segment of
0.1 cm (0.319 cm in Test Case 1), which sets the length scale: the
dimensionless parameters are $Q\simeq122.1$, $V\simeq6.38$ (printed as
\texttt{Phi}) and $F_{ve}\simeq122.1$ in the definition of Reneker et al.\
(1727.5 with the definition of Tab.~\ref{tab:tabella-adimensionale}). Here
we have activated the injection of new beads by the directive
\texttt{inserting yes} and their removal at the collector by
\texttt{removing yes}; the run uses the fourth-order Runge--Kutta integrator
(\texttt{integrator 3}) with a timestep of $10^{-7}$ s up to a final time of
2 s. In the OpenACC build this one-dimensional case runs the CPU build's
code, since one-dimensional systems do not take the device step; its
Coulomb sum is computed on the GPU, with copies at every call, only while
the jet has at least 128 active beads.
